\chapter{Non-Amplification and Capacity Bounds}

\section{Purpose of Non-Amplification}

The non-amplification principle records a simple structural rule: an
admissible refinement process cannot extract more relevant information than
its observation model permits. If the process has bounded locality, bounded
message size, bounded detector access, or bounded update width, then each step
has a finite information budget.

The principle is not a claim that information can never increase in any
physical or computational sense. It is a claim about admissible refinement
systems after the allowed observations and updates have been fixed. Once that
admissible class is fixed, the process cannot use information that is outside
its permitted view.

This chapter makes the capacity assumption explicit and separates three
levels:

\begin{enumerate}
\item local observation capacity;
\item one-step entropy reduction;
\item cumulative refinement depth.
\end{enumerate}

\section{Locality of Refinement Processes}

Let \(R\) denote a locality radius.

\begin{definition}[\(R\)-local refinement]
A refinement step is \(R\)-local if the update at a site \(v\) depends only on
data inside the radius-\(R\) neighborhood of \(v\).
\end{definition}

In a graph model, the neighborhood is the graph ball \(B_R(v)\). In a logical
model, locality may be expressed by bounded-variable or bounded-radius
formulas. In an operational physical model, locality may be expressed by
finite detector access or bounded propagation.

\begin{definition}[Local view]
The local view at \(v\) is the information available to the update rule at
that site. It may include labels, adjacency data, bounded measurements, or
other admissible local features.
\end{definition}

The local view is not automatically bounded. Boundedness must be proved from
the model assumptions, such as bounded degree, fixed radius, finite alphabet,
finite detector precision, or finite register width.

\section{Information Capacity Bound}

Capacity bounds the information removed by one admissible step.

\begin{definition}[Information capacity]
Let \(H\) be an entropy functional on a state space \(X\), and let
\(\mathcal R\) be a class of admissible refinements. The class
\(\mathcal R\) has capacity \(C\) if every \(R\in\mathcal R\) satisfies
\[
H(x)-H(Rx)\le C
\]
for every \(x\in X\).
\end{definition}

The capacity \(C\) may represent bits, local types, rank, detector outcomes,
or another information measure. What matters is not the name of the unit, but
the existence of a uniform upper bound on one-step entropy reduction.

\begin{example}[Finite local alphabet]
Suppose every local update returns one of \(B\) possible local types. Then a
single local observation can distinguish at most \(B\) alternatives and carries
at most
\[
\log B
\]
units of information.
\end{example}

\section{From Locality to Capacity}

Locality implies a capacity bound only when the local view is finite or
effectively bounded.

\begin{theorem}[Locality-to-Capacity Bound]
Assume an admissible refinement step is \(R\)-local, the local alphabet has
size at most \(A\), and each update consults at most \(M\) local entries. Then
the number of possible local views is at most \(A^M\), and the information
available to one update is at most
\[
M\log A.
\]
\end{theorem}

\begin{proof}
Each of the \(M\) consulted entries has at most \(A\) possible values. Hence
the total number of possible local views is at most \(A^M\). The information
contained in a choice among \(A^M\) possibilities is
\[
\log(A^M)=M\log A.
\]
\end{proof}

\begin{remark}
The theorem gives a capacity estimate for one local view. A full parallel
system still requires an accounting rule specifying whether updates are
counted per site, per round, per block, or per global refinement step.
\end{remark}

\section{Non-Amplification Theorem}

The non-amplification theorem is the cumulative form of the capacity bound.

\begin{theorem}[Non-Amplification]
Let \(x_0,x_1,\ldots,x_t\) be an admissible refinement trajectory. Suppose
each step satisfies
\[
H(x_i)-H(x_{i+1})\le C.
\]
Then
\[
H(x_0)-H(x_t)\le tC.
\]
\end{theorem}

\begin{proof}
Summing the one-step inequalities gives
\[
\sum_{i=0}^{t-1}\bigl(H(x_i)-H(x_{i+1})\bigr)\le \sum_{i=0}^{t-1}C=tC.
\]
The left side telescopes to \(H(x_0)-H(x_t)\).
\end{proof}

\begin{theorem}[Capacity Depth Bound]
If terminal resolution requires \(H(x_t)=0\), then every admissible trajectory
from \(x_0\) to a terminal state satisfies
\[
t\ge \frac{H(x_0)}{C}.
\]
\end{theorem}

\begin{proof}
By non-amplification,
\[
H(x_0)-H(x_t)\le tC.
\]
If \(H(x_t)=0\), then \(H(x_0)\le tC\), and therefore
\[
t\ge \frac{H(x_0)}{C}.
\]
\end{proof}

\section{Parallel Refinement}

Parallelism does not remove the need for capacity accounting. It changes the
unit being counted.

\begin{definition}[Parallel round]
A parallel round is a collection of local updates applied simultaneously
according to an admissible synchronization rule.
\end{definition}

A parallel round may remove more total information than a single local update.
However, the round still has a capacity bound if the number of independent
updates and the capacity of each update are bounded.

\begin{theorem}[Parallel Capacity Bound]
Suppose a parallel round consists of at most \(m\) independent local updates,
and each local update removes at most \(C_{\mathrm{loc}}\) units of entropy.
Then one parallel round removes at most
\[
mC_{\mathrm{loc}}
\]
units of entropy.
\end{theorem}

\begin{proof}
The entropy removed by the round is at most the sum of the entropy removed by
its local updates. Each local update removes at most \(C_{\mathrm{loc}}\), and
there are at most \(m\) updates. Therefore the total is at most
\(mC_{\mathrm{loc}}\).
\end{proof}

If \(m\) grows with the system size, the round capacity may grow with the
system size. In that case, a linear-depth conclusion does not follow from a
constant-capacity argument. The model must state whether capacity is measured
per update, per round, per processor, per unit area, or globally.

\section{Algorithmic Implications}

For algorithmic refinement systems, capacity bounds say that a local update
cannot distinguish exponentially many hidden alternatives unless its local
view also grows.

\begin{example}[Search space]
Let the unresolved search space have size \(2^n\), so that \(H_0=n\) bits. If
each admissible refinement step removes at most \(c\) bits, then terminal
resolution requires at least
\[
\frac{n}{c}
\]
steps.
\end{example}

This does not prove a lower bound for arbitrary algorithms. It proves a lower
bound only for the refinement class whose capacity has been bounded.

\section{Physical Implications}

For physical systems, capacity bounds express finite operational access. A
detector, channel, boundary register, or finite-energy measurement apparatus
can expose only a bounded amount of information per admissible operation.

\begin{example}[Finite detector outcomes]
Suppose a detector has \(B\) distinguishable outcomes per measurement. Then
one measurement carries at most \(\log B\) units of outcome information. If a
state has \(H_0\) units of unresolved operational entropy, at least
\(H_0/\log B\) such measurements are required to resolve it, assuming each
measurement is admissible and independent in the required sense.
\end{example}

The example is an operational accounting statement, not a complete physical
theorem. A physical theorem must also specify dynamics, measurement coupling,
noise, backreaction, and admissibility.

\section{Failure Modes}

A claimed non-amplification result fails if any of the following occur:

\begin{enumerate}
\item the update has hidden global access;
\item the local alphabet is unbounded;
\item the locality radius grows with system size;
\item the number of parallel updates is counted incorrectly;
\item the entropy functional is not monotone under refinement;
\item the terminal condition does not require the claimed entropy collapse.
\end{enumerate}

These are not exceptions to the theorem. They are failures of the hypotheses.

\section{Capacity Audit Checklist}

A capacity-bound argument should provide the following data:

\begin{center}
\begin{tabular}{ll}
\toprule
Field & Required content \\
\midrule
State space & configurations under refinement \\
Entropy & unresolved information or defect measure \\
Locality & radius, detector scope, or observation class \\
Alphabet & possible local observations or outcomes \\
Step rule & admissible refinement operation \\
Capacity & one-step entropy reduction bound \\
Parallel rule & accounting convention for simultaneous updates \\
Terminal state & condition defining resolution \\
\bottomrule
\end{tabular}
\end{center}

\section{Boundary}

This chapter proves only capacity bookkeeping for admissible refinement
systems. It does not prove that all computation is local, does not prove that
all physical systems have bounded operational capacity, and does not prove
unrestricted Chronos-RR, unrestricted H4.1/FGL, P vs NP, or any Clay problem.
